\chapter{Radar Clubhead Kinematics: Motion and Observability}
\label{app:screw}
% Retain the historical background reference without claiming a new full-book read.
\nocite{chenmicrodoppler}

A useful club-motion model must connect what the sensor records to what a golfer wants to know: which point on the head is moving, how the face is turning, and how confidently those quantities are known near impact.
Screw theory provides a compact representation of rigid motion, but choosing that representation does not make the motion observable.
This appendix develops the measurement equations, exposes an exact ambiguity in instantaneous single-origin Doppler data, and specifies what an implementer must add and validate before reporting club-delivery parameters.
The estimator described below is a research design, not a demonstrated IWR6843 golf instrument.
The point-tracking framework in \cref{app:ekf} is a starting point; a rigid-body estimator also needs orientation, a calibrated head reference, and an explicit account of unresolved motion.

\section{Context: Definitions Are Not an Estimator Specification}
\label{sec:screw-context}

TrackMan defines club speed at the head's geometric center immediately before first contact, while its path and attack-angle definitions use maximum compression as their time reference~\cite{trackmanclubspeed,trackman40params}.
These definitions specify intended outputs, not the complete sensing and estimation algorithm.
The manufacturer's discussion of a three-dimensional club silhouette and its patents describe additional information beyond one list of radial velocities; they do not establish that the particular Doppler-only solve below reconstructs a head pose~\cite{trackmanclubdata,us10850179}.
Nor does an algorithm's appearance in a patent establish its implementation or independently measured accuracy in every product version.
This appendix makes no claim about an undisclosed vendor state representation or the novelty or legal scope of the proposed design.

Instantaneous screw axes also have a distinct use in biomechanics.
Vena et al. studied optically reconstructed pelvis, shoulder, and left-arm motion: Part 1 reports predominantly rotational contributions to marker velocity in those segments, and Part 2 reports the expected segment sequence in two of five subjects~\cite{vena2010a,vena2010b}.
Those observations neither validate a radar clubhead estimator nor establish a universal axis near the hands.
The publisher abstracts and the visible Part 1 mathematical appendix were checked for this review; the subscription main texts were not accessed.
The robotics conventions below are independently supported by the authors' explanation of body and spatial twists~\cite{modernroboticstwists}; Murray, Li, and Sastry remain a general reference~\cite{murrayliss}.

\section{Rigid Motion With Declared Frames and Reference Points}
\label{sec:screw-primer}

Assume a rigid clubhead over the interval being modeled, a stationary radar origin, and a right-handed Cartesian sensor frame.
This is a head model: shaft bending and deformation during contact require separate states or an explicitly bounded approximation.
Let $O$ be a declared body-fixed reference point, $\vect{p}_O$ its sensor-frame position, $\vect{v}_O=\dot{\vect{p}}_O$ its velocity, and $\vect{\omega}$ the head's angular velocity expressed in the same frame.
Use linear-first ordering throughout:
\begin{equation}
\label{eq:screw-twist}
\xi_O=\begin{bmatrix}\vect{v}_O\\\vect{\omega}\end{bmatrix}\in\mathbb R^6.
\end{equation}
This six-component velocity description becomes a body or spatial Lie-algebra twist only after the reference and expression frame are specified below.
For a body-fixed point whose instantaneous sensor-frame offset from $O$ is $\vect r$,
\begin{equation}
\label{eq:screw-pointvel}
\vect v(\vect r)=\vect v_O+\vect\omega\times\vect r.
\end{equation}
Moving the reference to another body-fixed point $Q$ therefore changes the linear component to $\vect v_Q=\vect v_O+\vect\omega\times(\vect p_Q-\vect p_O)$.
The physical velocity field is unchanged; its six numerical coordinates are not invariant to changes of reference or frame.

For $\lVert\vect\omega\rVert>0$, the instantaneous screw axis is the line
\begin{equation}
\label{eq:screw-isa}
\begin{aligned}
\vect p_{\mathrm{axis}}(a)&=\vect p_O+\vect r_{\mathrm{ISA}}+a\uvec\omega,
&\vect r_{\mathrm{ISA}}&=\frac{\vect\omega\times\vect v_O}{\lVert\vect\omega\rVert^2},\\
h&=\frac{\vect\omega\cdot\vect v_O}{\lVert\vect\omega\rVert^2},
&\uvec\omega&=\frac{\vect\omega}{\lVert\vect\omega\rVert}.
\end{aligned}
\end{equation}
Here $a$ is distance along the axis and $h$ is signed axial translation per radian.
Substitution into \cref{eq:screw-pointvel} gives velocity $h\vect\omega$ at every point on the axis, which verifies the construction.
Pure translation has no unique finite axis, rest has no preferred axis, and near-zero angular speed makes the axis and pitch numerically sensitive.
An instantaneous axis is not a uniquely defined swing plane: its moving line need not remain normal to one plane, and a head-point trajectory need not be planar.
A plane estimate needs its own declared point trajectory, fitting interval, residual, and uncertainty.
Low pitch or an axis near the grip can be hypotheses to test; neither is an assumption that the cited segment study licenses for every club swing.

\subsection{Body and Spatial Twists for Pose Integration}

Let $T=\left[\begin{smallmatrix}R&\vect p_O\\0&1\end{smallmatrix}\right]$ map body coordinates to the sensor frame, with $R\in SO(3)$.
For a linear-first pair $V=(\vect v,\vect\omega)$, define $V^\wedge=\left[\begin{smallmatrix}[\vect\omega]_\times&\vect v\\0&0\end{smallmatrix}\right]$, where $[\vect a]_\times\vect b=\vect a\times\vect b$.
Then
\begin{equation}
\label{eq:screw-frames}
\begin{aligned}
V_b&=(R^{\mathsf T}\vect v_O,\ R^{\mathsf T}\vect\omega),
&V_b^\wedge&=T^{-1}\dot T,\\
V_s&=(\vect v_O-\vect\omega\times\vect p_O,\ \vect\omega),
&V_s^\wedge&=\dot T T^{-1}.
\end{aligned}
\end{equation}
These are related by the adjoint of $T$~\cite{modernroboticstwists}.
A constant body-twist step uses $T_{k+1}=T_k\exp(V_{b,k}^\wedge\Delta t)$; a constant spatial-twist step uses $T_{k+1}=\exp(V_{s,k}^\wedge\Delta t)T_k$.
For varying motion these are integration approximations; the increments must be composed in time order.
Inserting sensor-frame $\vect v_O$ directly into the right-multiplying body update generally gives the wrong trajectory.
Likewise, a logarithm of two measured poses describes a finite interframe displacement; dividing it by the interval is not automatically an unaliased instantaneous angular velocity.

\section{The Measurement Model and Its Exact Ambiguity}
\label{sec:screw-measurement}

First consider an ideal monostatic measurement represented by a single origin: a localized scattering point $\vect p_i\ne0$ has sight line $\uvec u_i=\vect p_i/\lVert\vect p_i\rVert$, offset $\vect r_i=\vect p_i-\vect p_O$, and radial velocity $z_i$ positive away from the radar.
Assume its Doppler return represents the instantaneous velocity of head material there.
The conditional measurement equation is
\begin{equation}
\label{eq:screw-doppler}
z_i=\uvec u_i^{\mathsf T}\vect v_O
 +(\vect r_i\times\uvec u_i)^{\mathsf T}\vect\omega+\epsilon_i.
\end{equation}
With the ordering in \cref{eq:screw-twist},
\begin{equation}
\label{eq:screw-ls}
\vect z=A\xi_O+\vect\epsilon,\qquad
A_i=\bigl[\uvec u_i^{\mathsf T}\mid(\vect r_i\times\uvec u_i)^{\mathsf T}\bigr].
\end{equation}
Linearity holds conditional on the positions, associations, and rigid-motion model.
It does not say that six independent unknowns can be recovered.
Because $\vect p_i\times\uvec u_i=0$, the same equation factors as
\begin{equation}
\label{eq:screw-rank}
\begin{aligned}
\vect z&=U\bigl(\vect v_O-\vect\omega\times\vect p_O\bigr)+\vect\epsilon,\\
A&=U\bigl[I\mid[\vect p_O]_\times\bigr],\qquad
\operatorname{rank}(A)\le3,
\end{aligned}
\end{equation}
where row $i$ of $U$ is $\uvec u_i^{\mathsf T}$.
For any $\delta\vect\omega$, the perturbation $\delta\vect v_O=\delta\vect\omega\times\vect p_O$ leaves every predicted Doppler value unchanged.
This is an exact three-dimensional null family even with arbitrarily many noncoplanar localized detections and noise-free measurements.
Nearly parallel sight lines can further weaken the remaining three combinations; they do not explain away the structural nullspace.
The observed combination is the translational component of $V_s$, not necessarily the speed of the head center.

\begin{keypoint}
One origin's instantaneous Doppler rows cannot determine the complete club velocity field under this model.
Better signal-to-noise ratio or more rows at that same origin cannot remove this ambiguity.
Range and angle histories, known feature geometry, independent orientation measurements, or sufficiently informative additional sensor geometry can supply information outside these rows; their contribution must appear in the actual measurement model.
\end{keypoint}

For example, known and correctly associated noncollinear body features with resolved three-dimensional positions can constrain pose; their time history can then constrain motion.
Unlabeled specular detections are not automatically such features.
Two separated monostatic stations can generically raise the Doppler-only rank to five, but retain rotation about their baseline as a null direction with its corresponding translational change.
Three noncollinear station origins can generically reach six with sufficient point/sight-line diversity.
These are geometric possibilities, not guarantees of adequate precision.
A multiple-input multiple-output (MIMO) array also requires its actual transmit/receive path model: a bistatic path derivative involves both ray directions.
Twelve virtual array channels are not twelve independent monostatic stations, and the single-origin bound is not a theorem that every near-field, coherent, or feature-tracking radar architecture is limited to three observable motion coordinates.

\subsection{What Least Squares and Covariance Actually Provide}

For fixed $A$ and zero-mean measurement error with known positive-definite covariance $R_z$, generalized least squares minimizes $(\vect z-A\xi)^{\mathsf T}R_z^{-1}(\vect z-A\xi)$.
If $A$ has full column rank, its estimator covariance is $(A^{\mathsf T}R_z^{-1}A)^{-1}$ under that model.
This inverse does not exist for the six-unknown single-origin problem.
A pseudoinverse selects a convention for the unobserved components; a zero eigenvalue in a pseudoinverse covariance is not evidence of zero uncertainty there.
Report the identifiable combinations and the remaining nullspace, or explicitly state the extra observations or priors that define a full posterior.

Inspect singular values after noise whitening and declared parameter scaling, for example $R_z^{-1/2}AD$ with $D$ converting dimensionless unknowns into chosen linear- and angular-velocity scales.
Raw singular values mix meters per second with radians per second, so an unexplained threshold is not a physical observability criterion.
Uncertain positions make $A$ uncertain; position--Doppler correlations, calibration bias, ambiguous associations, and model discrepancy require joint propagation or an appropriate errors-in-variables model.
Robust losses such as Huber or Cauchy generally require iterative fitting and their own uncertainty treatment.
They do not turn an incorrect scattering model into a calibrated rigid-body measurement.

\section{What Representative Hardware Adds}
\label{sec:screw-ceiling}

\subsection{OPS243-A: Radial-Speed Information}

The manufacturer's OPS243 product brief distinguishes the Doppler-only A model from the range-capable C model~\cite{ops243brief}.
Without per-return range and direction, the A output alone cannot assemble the located rows above.
A Doppler spectrum is a scattering-amplitude-weighted signal, not a uniform sample of clubhead material velocities.
Its peak need not always be the toe, its midpoint need not be the geometric-center speed, and its width is not a uniquely calibrated angular speed.
For a stated example of 30\,ksps, 128 samples, and 24\,GHz, $\Delta v=cf_s/(2f_cN)\approx1.46$\,m/s, or 3.27\mph{}, before accounting for the spectral window.
Zero-padding samples the spectrum more finely but does not add the observation time needed to resolve closer velocities.
A percentile or spectral-width feature may be useful after empirical calibration across aspect, club, speed, and clutter; it should retain that provenance.

\subsection{K-LD7: Geometry and Operating Envelope}

The K-LD7 Revision B datasheet specifies range, radial speed, and one angular coordinate, with a maximum speed setting of 100\,km/h and a 29\,ms cycle at that setting~\cite{kld7}.
The speed limit is approximately 27.8\,m/s, not a limit on total three-dimensional speed independent of viewing direction.
A head moving at 50\,m/s travels 1.45\,m during that cycle; its radial component may alias even though another orientation does not.
A single angular coordinate also does not localize arbitrary three-dimensional scatterers.
Neither these specifications nor a frame's angle output establishes that high-speed ball returns remain valid angle measurements after Doppler aliasing.
Qualification would need the chosen geometry, waveform, association, and motion envelope.

\subsection{IWR6843: A Coupled Waveform Budget}

The IWR6843 has three transmitters, four receivers, a 60--64\,GHz band, and a maximum specified ramp rate of 250\,MHz/\si{\micro\second} in Revision F~\cite{iwr6843}.
Array angle estimation and multiplexing depend on antenna geometry and processing; they are not generically synonymous with monopulse~\cite{timimo}.
Let $T_c$ be the start-to-start chirp interval, $n_T$ the number of cyclic time-division-multiplexed transmitters, and $N$ the number of repeated samples \emph{per transmitter}.
For uniform slow-time sampling, wavelength $\lambda$, and a conventional unambiguous Doppler FFT,
\begin{equation}
\label{eq:screw-chirp}
T_r=n_TT_c,\qquad |v_r|<\frac{\lambda}{4T_r},\qquad
\Delta v=\frac{\lambda}{2NT_r},\qquad T_{\mathrm{obs}}\approx NT_r.
\end{equation}
These follow from the round-trip phase increment $4\pi v_rT_r/\lambda$ and the sampled Fourier frequency grid~\cite{tichirp,timimo}.
At the illustrative $\lambda=5$\,mm, $T_c=25$\,\si{\micro\second}, and $N=128$:
\begin{center}\small
\begin{tabular}{@{}lrrr@{}}
\toprule
\textbf{Transmit Schedule} & \textbf{Speed Limit (m/s)} & $\boldsymbol{\Delta v}$ \textbf{(m/s)} & $\boldsymbol{T_{\mathrm{obs}}}$ \textbf{(ms)} \\
\midrule
One transmitter & 50.0 & 0.781 & 3.2 \\
Three-transmitter TDM & 16.7 & 0.260 & 9.6 \\
\bottomrule
\end{tabular}
\end{center}
The table is arithmetic, not a validated firmware configuration; frame idle time and processing can lower the output rate further.
At the sampling boundary the signed velocity is ambiguous, so a 50\,m/s limit provides no margin for a 50\,m/s return.
For three-transmitter TDM to exceed that unambiguous radial speed at this wavelength requires $T_c<8.33$\,\si{\micro\second}.
Sweeping 4\,GHz already needs at least 16\,\si{\micro\second} at the stated maximum ramp rate, before settling and idle time.
Thus those requirements cannot be combined in that simple waveform; ambiguity-resolution schemes or other multiplexing designs need their own analysis and validation.

Nominal range resolution is $c/(2B)\approx3.75$\,cm for an effective sampled bandwidth $B=4$\,GHz~\cite{tichirp}.
It does not guarantee three or four distinct returns from a head: projected radial extent, reflectivity, angular resolution, windowing, and detection thresholds also matter.
Even over 3.2\,ms, a 50\,m/s radial return migrates 16\,cm, more than four such bins.
Acceleration, changing aspect, range migration, and inter-transmitter phase changes must therefore be included in the signal model or shown to be negligible for a bounded operating case.
TI specifically describes motion-phase correction before TDM angle estimation~\cite{timimo}.
Select sampled bandwidth, slope, ADC start/sample time, IF bandwidth, settling, transmitter schedule, output timing, and processing load together~\cite{tichirp}; a bandwidth, frame rate, and velocity ceiling cannot be independently promised.

\section{An Estimation Recipe With Explicit Acceptance Conditions}
\label{sec:screw-recipe}

\subsection*{Step 1 --- Preserve Geometry, Timing, and Detection Uncertainty}

Retain the raw signal or sufficient diagnostics to check aliasing, motion compensation, and false detections.
Choose CFAR training/guard cells and any order-statistic variant against recorded clutter and false-alarm requirements; no detector is universally preferable for every impact scene.
For a declared right-handed sensor frame with range $\rho$, azimuth $\theta$, and elevation $\phi$, convert spherical coordinates to Cartesian positions:
\begin{equation}
\vect p=\rho\begin{bmatrix}\cos\phi\cos\theta&\cos\phi\sin\theta&\sin\phi\end{bmatrix}^{\mathsf T}.
\end{equation}
Propagate the full range/angle covariance through this mapping and preserve relevant correlations with Doppler.
A reported FFT bin spacing, a precision estimate, and demonstrated accuracy are different quantities.
Assign timestamps to the acquisition intervals; a frame output timestamp is not an instantaneous sample of all returns.

\subsection*{Step 2 --- Associate Returns With a Head Model}

Use range, radial velocity, geometry, temporal continuity, and calibrated uncertainties jointly to distinguish head, shaft, ball, and body returns.
A shaft point is not invariably slower in radial speed or nearer the sensor than the head, and speed alone cannot guarantee removal of body clutter.
An SNR-weighted detection centroid can migrate as aspect changes; its derivative includes that migration and is not automatically the velocity of a body-fixed point.
Choose $O$ through a calibrated geometric model or identified body feature, and distinguish material-point velocity from apparent glint motion.
If correspondence or head geometry is unresolved, retain multiple hypotheses or withhold outputs that depend on it.

\subsection*{Step 3 --- Fit the Identifiable Quantities First}

Use \cref{eq:screw-rank} to fit the three-component combination $\vect b=\vect v_O-\vect\omega\times\vect p_O$ when the weighted sight-line matrix supports it.
A smaller observable subspace may be all the data justify.
Record residuals, scaled singular values, associations, and the acquisition interval.
Do not publish an arbitrary six-component pseudoinverse solution as a fully measured twist.
A full solve needs additional, explicitly modeled observations, such as calibrated feature-position histories or independent orientation information.

\subsection*{Step 4 --- Establish Temporal Observability}

A changing sight line is not a proof that smoothing recovers six arbitrary motion functions.
For Doppler-only histories, even a constant rotation of every head position and velocity about the sensor preserves all radial histories while changing the target-relative motion; both trajectories can be smooth.
Angle or feature-position observations may distinguish them, which is precisely why those observations must be included rather than used only as fixed coefficients in a velocity solve.
Analyze the linearized multi-time measurement model with the chosen state transitions, geometry, calibration states, and priors.
Report structural nullspaces separately from poor conditioning and from uncertainties controlled chiefly by a prior.
Use withheld trajectories and deliberate near-degeneracies to test the claimed operating region.

\subsection*{Step 5 --- Estimate Pose and Motion Consistently}

An $SE(3)$ filter or batch smoother can combine pose, velocity, calibration, and feature observations using \cref{eq:screw-frames}.
Feed raw residuals where practical, so the estimator retains their actual information and correlations; a per-frame diagnostic fit can still be useful.
A backward smoother uses future measurements and must be labeled offline if they are unavailable at the desired reporting time.
Select process uncertainty against observed dynamics, not an unsupported universal acceleration of one mile per hour per millisecond.
Low-pitch or grip-proximity penalties are optional model assumptions whose sensitivity must be reported; they cannot resolve an unobservable quantity through measurement alone and become singular near zero angular speed if written using \cref{eq:screw-isa} without safeguards.
A prior that is weak in observed directions may completely determine an unobserved one.

\subsection*{Step 6 --- Define the Impact Event and Model Boundary}

Synchronize the radar, any optical reference, and the contact-event measurement, including sound-propagation and trigger latency where an acoustic trigger is used.
Distinguish the last resolved pre-contact state, first contact, maximum compression, and separation.
A prediction from the last pre-contact sample to another event is an extrapolated estimate with event-time and model uncertainty.
The manufacturer's historical club-data explanation explicitly discusses using pre-impact data for an event-time definition~\cite{trackmanclubdata}; the chosen product definition must accompany any comparison.
Do not assume that every head is exactly rigid until a universal 500\,\si{\micro\second} switch, or extend a pre-contact process model through collision and turf contact without validation.
Fit pre- and post-contact intervals separately if the purpose is to study collision-induced changes.

\subsection*{Step 7 --- Report Only Supported Delivery Quantities}

Let the target frame be right-handed with $x$ down the target line, $y$ to the golfer's left when looking down that line, and $z$ upward.
For a calibrated body-fixed point with body offset $\vect r_{gc}^{,b}$, first compute $\vect v_{gc}=\vect v_O+\vect\omega\times R\vect r_{gc}^{,b}$ and express it in that target frame.
Then speed is $\lVert\vect v_{gc}\rVert$, path positive to the right is $\operatorname{atan2}(-v_y,v_x)$, and attack angle positive upward is $\operatorname{atan2}(v_z,\sqrt{v_x^2+v_y^2})$.
Path is undefined for zero horizontal velocity, and the velocity direction is undefined at rest.
A calibrated geometric-center offset, an anatomical grip point, and a changing radar centroid are not interchangeable.

Face closure needs a definition and orientation information.
For a known target-frame unit face normal $\vect n$, rigid rotation gives $\dot{\vect n}=\vect\omega\times\vect n$.
Its mathematical horizontal azimuth $\gamma=\operatorname{atan2}(n_y,n_x)$ satisfies
\begin{equation}
\label{eq:screw-closure}
\dot\gamma=\frac{n_x\dot n_y-n_y\dot n_x}{n_x^2+n_y^2},
\qquad n_x^2+n_y^2>0.
\end{equation}
This generally differs from $\vect\omega$ projected onto the shaft axis.
A right-positive face-angle convention reverses this azimuth sign; a label such as closing also needs the handedness and stated convention.
A full head frame requires more orientation information than one normal, and face angle at the contact location may additionally require head geometry and impact location.
Thus closure rate is not a free Doppler-only output or a parameter whose product availability can be declared universally exclusive.

Estimate swing plane and swing direction only through a specified trajectory-plane fit, with its residual and fit window.
A low point requires an observed or modeled trajectory with $v_z=0$ and a local height minimum; a vertical-velocity zero alone can be a maximum.
A difference between pre- and post-impact angular velocities is a net kinematic change.
Attributing it to gear effect requires an impact/contact model, timing, inertia, and consideration of continuing shaft/grip and other contact loads.

Propagate joint state, calibration, association, timing, and model uncertainty into each output and disclose unresolved directions or dependence on priors.
Keep the book's distinction between sensor observations, quantities derived from them, and model-based estimates; merely changing the state representation promotes no output into a measured category.
Kinematics alone also identifies neither golfer torques nor a control strategy: a hand-axis prior must not be presented as evidence that the golfer actually produced that motion by the assumed mechanism.

\section{Validation That Can Falsify the Proposed Estimator}
\label{sec:screw-validation}

The following is a proposed qualification sequence, not a report of completed hardware or human testing.
Use the release and evidence discipline of \cref{ch:validation-program}.
\begin{enumerate}
\item \textbf{Analytic and Synthetic Checks.}
Verify the exact null family, frame conversions, pure-translation limit, and waveform arithmetic before fitting realistic data.
Simulations should include noise, glint changes, aliasing, missed detections, correlated errors, and deliberately unobservable cases, and should test empirical interval coverage as well as point error.
Recovering truth in a simulation generated with the same restrictive prior cannot establish that the prior holds for golfers.
\item \textbf{Independent Rig Reference.}
Use synchronized calibrated encoders or optical tracking for pose and velocity on a rigid-motion rig.
A fixed-axis pendulum supplies a known axis only if the bearings, mounting, head rigidity, and coordinate calibration support that model; its angular history is not known merely because it is called a pendulum.
Test translations, multiple axes, low angular speeds, varying aspects, and moving-pivot cases as well as the easy planar case.
\item \textbf{Independent Cross-System Comparison.}
Match the head point, event time, coordinate convention, club, environment, and uncertainty of the reference.
A commercial monitor is a comparator, not automatically ground truth; unavailable or model-derived comparator outputs cannot validate an independent direct measurement.
Report bias, limits of agreement, repeatability, missing-output rates, and uncertainty coverage against predefined requirements, including difficult cases.
\item \textbf{Mechanics Consistency With Stated Assumptions.}
Check geometric identities only where their planar or rigid-body assumptions apply.
An identity satisfied by construction is a software consistency check, not independent physical validation.
Do not discard shots merely because a simplified smash-factor or swing-plane model is violated: inspect reference-point mismatch, contact, alignment, and model inadequacy, and preserve the rejection rule and excluded cases.
\end{enumerate}

\begin{implication}
The value of a screw representation is a common account of rigid motion and point transport.
Its practical benefit depends on calibrated geometry, sufficient observations, consistent timing, and validated uncertainty.
Declare what the data constrain before reporting a complete club state, and keep the inferred motion separate from claims about the forces or intentions that produced the swing.
\end{implication}
